problem:  
Let \((M^{2n+1}, \alpha, \omega, J, g)\) be a compact, oriented smooth manifold equipped with a **quasi-contact metric structure** in the following refined sense:  
- \(\alpha\) is a unit \(1\)-form defining the cooriented distribution \(\xi = \ker \alpha\);  
- \(\omega\) is a \(2\)-form such that \(d\omega = 0\) and \(\alpha \wedge \omega^n > 0\) (so \(\omega\) is symplectic on each even-dimensional leafwise component of \(\xi\), but \(d\alpha|_{\xi}\) may be degenerate);  
- \(J:\xi \to \xi\) is an almost complex structure with \(\omega(X,Y)=g(JX,Y)\) for \(X,Y \in \xi\); and \(R\) is the Reeb-like vector field uniquely determined by \(\alpha(R)=1,\ \iota_R\omega=0\).  
Assume that \(R\) is geodesic and that \(g\) has **positive mixed sectional curvature**:
\[
\sec^g(X,R) > 0 \quad \text{for all } 0 \ne X \in \xi.
\]

Now consider the space of deformations \((\alpha_t, \omega_t)\), \(t \in [0,1]\), such that  
- \(\alpha_t \wedge \omega_t^n > 0\) and \(d\omega_t = 0\) for all \(t\);  
- \(\alpha_0 = \alpha\), \(\omega_0 = \omega\), and \((\alpha_1,\omega_1)\) defines a **contact metric structure**, i.e. \(\omega_1 = d\alpha_1\) and \(d\alpha_1|_\xi\) is nondegenerate.  

**Problem:**  
If \(g\) has positive mixed sectional curvature and \(R\) integrates to an isometric circle action, does there always exist a finite covering \(\tilde M \to M\) such that the pulled-back structure \((\tilde \alpha, \tilde \omega, \tilde g)\) admits a deformation of the above type to a contact metric structure \((\tilde \alpha_1, \tilde \omega_1, \tilde g_1)\)?  

Equivalently: does positive mixed curvature combined with an isometric Reeb flow enforce the “closure” of a quasi-contact structure into a contact one after passing to some finite cover?

Why is it a "good" problem:  
This problem carefully intertwines **odd-dimensional curvature geometry** with the **topological rigidity** of quasi-contact structures, within a precisely defined setting compatible with standard frameworks (almost contact metric or stable Hamiltonian structures). The assumption that the Reeb-like flow is isometric places the problem at the interface of **Riemannian foliations** and **contact/Sasakian geometry**, where curvature restrictions are known to have profound topological implications.  

Its simplicity of statement belies its depth: the question asks whether a purely **metric positivity condition**—positive mixed sectional curvature—forces a topological and geometric transition from a quasi-contact form (where \(d\alpha|_{\xi}\) may be degenerate) to a fully contact form, possibly after finite cover. Resolving this could uncover a new kind of **Riemannian-to-topological rigidity**, extending principles like Myers' and Synge's theorems into the domain of **contact-like geometries**.  

It connects multiple areas—**contact topology**, **Riemannian curvature**, and **geometric dynamics of Reeb flows**—and may lead to new invariants or obstructions linking curvature and contactness. Its formulation is mathematically sound, its outcome unknown, and it has the potential to open a new direction of research, satisfying the core criteria for a “good” problem.